\section{Recommendations}
\label{sec:recommendations}

The project guidelines pose a simple closing question: what should be
done differently based on the evidence? The answer is grouped into
four audiences, each with a single, narrowly scoped recommendation
and the specific number from the analysis that justifies it.

\subsection{For RBI liquidity managers: monitor the WACMR-Repo
spread as a real-time barometer}

\textbf{Evidence.} The engineered \texttt{spread\_wacmr\_minus\_repo}
is the fourth-ranked SHAP feature, with mean absolute SHAP 0.0725, a
larger predictive weight than the MSF Rate itself (0.0585). In
Regime 1 (Accommodation), the average spread is $-0.19$ percentage
points; persistent positive spreads above $+0.10$ percentage points
have historically preceded liquidity-deficit interventions.

\textbf{Recommendation.} Operationalise a real-time WACMR-Repo
spread dashboard inside the RBI Financial Markets Operations
Department, with a soft alert when the spread exceeds $+0.10$
percentage points for three consecutive weeks and a hard alert at
$+0.25$ percentage points. The model treats this single engineered
spread as more informative than any individual corridor rate. A
dashboard is cheap to build (the underlying data is already public)
and would surface emerging liquidity stress earlier than the current
monthly review cycle.

\subsection{For institutional bond and money-market desks: size
positions to model regime confidence}

\textbf{Evidence.} The walk-forward XGBoost achieves 70.9\%
directional accuracy and RMSE 0.1019 across 389 out-of-sample weeks.
The regime-conditional SHAP analysis shows that absolute SHAP
magnitudes are roughly twice as large in the tightening regime as in
the accommodation regime, which means the model is materially more
confident inside Regime 1 than inside Regime 0.

\textbf{Recommendation.} Use the model's directional forecast as a
trade signal, but size positions proportional to the cluster
distance ratio (\texttt{cluster\_dist\_0 / cluster\_dist\_1}). Near
the regime boundary, where the ratio is close to one, reduce
position sizing and widen bid-offer spreads. Deep inside either
regime, scale up to the desk risk limit. This is a simple,
implementable adjustment that uses the model's own implicit
confidence rather than treating every prediction as equally reliable.

\subsection{For equity and forex desk practitioners: stop using
equity and currency momentum as a WACMR input}

\textbf{Evidence.} Adding twenty-eight equity and currency features
(Nifty 50 OHLCV, USD/INR OHLCV, and ten custom technical indicators
across both instruments) did not move any feature into the SHAP top
15. The first market feature to appear is at rank 23 with mean
absolute SHAP 0.0038, two orders of magnitude smaller than the
corridor rates. This is a null result, but it is a strong null
result: the model had every opportunity to learn from these features
and chose not to.

\textbf{Recommendation.} Practitioners who maintain technical
indicators on Nifty 50 or USD/INR for the purpose of forecasting
overnight call-money rates should reallocate that screen real estate.
Whatever transmission mechanism connects equity and currency markets
to overnight rupee liquidity operates at lower frequency than the
weekly signal studied here, so monthly or quarterly reviews of these
relationships make more sense than continuous monitoring.

\subsection{For NDAP and the open-data community: publish T-Bill
auction cut-off yields as a dedicated weekly series}

\textbf{Evidence.} The \texttt{Treasury\_Bills\_Details.csv} table on
NDAP labels column \texttt{I7504\_10} as a yield metric, but the
actual values in the downloaded CSV are outstanding amounts in INR
crore. The 91-day, 182-day, and 364-day cut-off yields are widely
referenced in RBI publications but are not available as a
machine-readable weekly series through the open API.

\textbf{Recommendation.} Adding a confirmed, API-accessible weekly
yield series for the three T-Bill tenors is the single highest-value
data addition NDAP could make for fixed-income research on Indian
markets. The same recommendation applies to intraday LAF utilisation
data, which would let researchers separate intra-week stress events
that the weekly aggregation in this project necessarily masks.

\subsection{For policy communication and education: open-data
dashboards demonstrably work}

\textbf{Evidence.} The interactive system documented in
Section~\ref{sec:system} (a public dashboard, a counterfactual
simulator, and an LLM-powered data agent) is built entirely on open
NDAP and Yahoo Finance data, runs on free-tier infrastructure
(Vercel and Render), and serves all 545 weeks of the panel through a
single public URL. Production performance metrics
(Section~\ref{sec:qa}) show sub-second median page load times across
ten pages on both desktop and mobile.

\textbf{Recommendation.} Public-facing tools that explain the
mechanics of monetary-policy transmission do not require private data
or expensive infrastructure. The implication for RBI's communications
team and for academic researchers is that lightweight open-data
dashboards can be deployed within a day or two of analytical work and
provide a far more direct way to communicate quantitative findings to
non-specialist audiences than a written research note alone.
